\chapter{Creazione Model}\label{chap:Model}
\section{Mapping delle classi}
Il model dell'applicativo è suddiviso in parte nel package ``DatabaseManagementPackage'' e in parte, per
le parti non legate al database, nel package principale.
Il package di interfaccia al database è stato realizzato partendo, mediante \ac{EF Core},
dalle tabelle del database, generando le classi entity.
Una volta generate le classi entity, si sono individuati gli oggetti fondamentali del dominio,
che sono stati mappati in classi \ac{DTO} mediante il package ``Mapster''.
Questi oggetti sono stati scelti in quanto rappresentano le entità principali del dominio,
come ad esempio le entità ``Cliente'', ``Prodotto'' e ``Documento''.
Le classi \ac{DTO} sono state realizzate partendo dalla struttura delle entity e, in
corrispondenza dei campi delle entity legati alle foreign key, sono stati adottati due approcci differenti.
Se il campo aveva una corrispondenza con una \ac{DTO}, si richiamava quest'ultima;
altrimenti si è deciso di mappare il campo \textit{Dictionary<int,string>}, dove si salvavano l'ID e la descrizione.
Il mapping è bidirezionale per poter riconvertire, quando necessario, un \ac{DTO} in entity per poterlo inserire nel database.
\section{Struttura del package}
Il package è stato strutturato in modo da separare le classi di test, i repository e le classi e interfacce esposte.
All'esterno del package oltre alle classi \ac{DTO} ed Enum relativi ad alcuni Lookup, è
stata esposta una classe univoca che espone i vari repository, denominata ``DatabaseManager''.
Questa soluzione permette di avere un unico punto di accesso ai repository, 
evitando di esporre le classi interne del package.
Il vantaggio della classe singola si ha anche nella gestione della Dependency Injection, 
in quanto si può registrare un'unica classe
per tutti i repository, evitando di registrare ogni singolo repository.
Si è poi realizzata una classe di test per ogni repository, che ne testasse le funzionalità di base.
Nel package è stata inoltre predisposta una classe con compito di mantenere salvati i dati dell'utente loggato,
come ad esempio l'ID dell'utente, il nome e cognome, il ruolo e i permessi.
Questa classe, registrata come singleton, permette di avere un unico punto di accesso ai dati dell'utente loggato,
permettendo la gestione dei permessi in ogni parte dell'applicativo.
\subsection{Classi di Test}
Le classi di test realizzate sono state pensate con l'obiettivo di funzionare anche in presenza di dati reali.
Per realizzare ciò si sono creati prima di ogni test dei savepoint del database,
che permettessero di tornare allo stato iniziale.
Successivamente, per ogni test, si creavano i dati di test fittizi che potessero validare
le funzionalità del repository senza
il rischio di prelevare dati reali.
Questo approccio ha richiesto molto tempo per la ricostruzione dei test più complessi, ma
ha permesso di testare piccoli cambiamenti senza dover ripristinare il database.
A fronte del tempo impiegato prima si è utilizzata \ac{AI} generativa per creare le bozze dei test,
 che successivamente sono state modificate e, in alcuni casi, riscritte completamente,
per adattarle al contesto reale.
Per realizzare ciò si è adottato un approccio articolato come segue:
\begin{enumerate}
    \item Si sono scritti i metodi del repository.
    \item Si sono generati i test sulla base dei metodi del repository.
    \item Si sono corretti i test per renderli accurati e validi.
    \item Si sono corretti i metodi del repository sulla base del risultato dei test.
\end{enumerate}
Questo approccio ha permesso di avere un ciclo di sviluppo più veloce,
 riducendo il tempo di scrittura dei test.
\subsection{Repository}
I repository sono stati divisi in base alle entità principali del dominio, come ad esempio i Prodotti.
Ogni repository poteva a sua volta gestire più \ac{DTO}, come ad esempio \textit{ProdottoDTO}
e \textit{ProdottocomposizioneDTO} appartenenti allo stesso oggetto di dominio.
Oltre alle operazioni \ac{CRUD}, si è scelto per quanto riguarda le operazioni di ricerca, di utilizzare
dei metodi che permettessero di inserire molteplici criteri di ricerca.
All'interno sono poi stati utilizzati i Queryable per costruire la query in esame.
\lstinputlisting[
  language=CSharp,
  caption={Esempio di query Queryable},
  label={lst:esempio-query-queryable}
]{listings/Queryable.cs}
Questo approccio ha permesso di costruire query complesse, senza dover scrivere molteplici metodi di ricerca.
Un'altra scelta che ha permesso di ridurre il codice ripetitivo
è stata quella di rimappare le nuove entità, create in \ac{DTO} dalle classi
di model del package principale, in classi di entity prive di Id che fossero pronte ad essere inserite nel database.
Tutti i metodi all'interno dei repository sono stati realizzati come asincroni,
per evitare di bloccare il thread principale dell'applicativo e, in particolare, l'interfaccia grafica.
\section{Gestione dei dati e delle transazioni}
Si è deciso di gestire le transazioni in modo atomico, considerandole come commit automatici,
per evitare una perdita dei dati in caso di errori o eccezioni.
I dati, prima dell'effettivo inserimento nel database,
vengono gestiti in memoria e, solo al termine di tutte le operazioni, vengono salvati nel database.
Questo approccio permette di non impegnare il database con dati parziali,
ma, al contempo, una volta validati, la transazione si può considerare come avvenuta, senza dover gestire rollback o commit.
All'interno dei repository, è stato aggiunto ai parametri passati ai metodi il ``CancellationToken''\cite{CancToken},
che permette di interrompere l'esecuzione di una query in caso di timeout o di chiusura dell'applicativo.
L'utilizzo del ``CancellationToken'' si può notare anche nella query di esempio~\ref{lst:esempio-query-linq},
 dove viene passato come parametro al metodo \texttt{ToListAsync()}.
Questa scelta è stata fatta per evitare di bloccare il database con query che non terminano,
e per permettere all'utente di cambiare schermata anche se la query non è terminata,
evitando di bloccare l'interfaccia grafica di destinazione.
